\color{unidarkgreen}{\subsection{Контроль цепей напряжения стороны НН}}
\color{black}

\begin{enumerate}

\item
Контроль цепей напряжения предназначен для выявления нарушений в цепях напряжения низкой стороны.

\item
Уставки КЦН представлены в таблице \ref{app_set:kcn}. На рисунке \ref{pic_kcn:KCN} приведен алгоритм работы КЦН. 

\begin{figure}[H]
\centering
\includegraphics[width=1\textwidth,height=1\textheight,keepaspectratio]{\fbpath/07.02 РЗАТ/041.1741 КЦН/img/KCN.pdf}
\captionof{figure}{Функциональная схема логики КЦН НН}\label{pic_kcn:KCN}
\end{figure}

\item
Контроль цепей напряжения низкой стороны (КЦН НН) основан на выявлении нарушения симметрии в цепях напряжения и просадки линейного напряжения. Порог срабатывания пуска по минимальному линейному напряжению задается уставкой \textbf{«Uл КЦН НН»}, порог срабатывания пуска по максимальному напряжению обратной последовательности задается уставкой \textbf{«U2 КЦН НН»}. Ввод в работу КЦН НН осуществляется с помощью программного переключателя \textbf{«КЦН НН»}. Выдержка времени контроля неисправности цепей напряжения регулируется уставкой \textbf{«Тср КЦН НН»}.

\item
Оперативный ввод/вывод КЦН НН производится с помощью БУ (блоков управления функциями) \textit{«Вывод КЦН НН»}. 


\end{enumerate}